\ExplSyntaxOn
\NewExpandableDocumentCommand{\Lang}{m m}
  { \bool_if:NTF \g_script_english_bool {#2} {#1} }
\NewExpandableDocumentCommand{\IfClozeTF}{m m}
  { \bool_if:NTF \g_script_cloze_bool {#1} {#2} }
\NewExpandableDocumentCommand{\IfEnglishTF}{m m}
  { \bool_if:NTF \g_script_english_bool {#1} {#2} }
\ExplSyntaxOff

\NewDocumentCommand{\EnglishTerm}{m}{%
  {\color{black!55}\itshape(engl.: #1)}%
}
\NewDocumentCommand{\GermanTerm}{m}{%
  {\color{black!55}\itshape(German: #1)}%
}

\NewDocumentCommand{\Cloze}{m m}{%
  \IfClozeTF
    {{\setlength{\fboxsep}{0pt}%
      \colorbox{black!12}{\makebox[#1][c]{\strut}}}}
    {#2}%
}

\NewDocumentCommand{\FrontMatterChapter}{m}{%
  \chapter*{#1}%
  \addcontentsline{toc}{chapter}{#1}%
}

% Grid dimensions are specified as numbers of 5 mm squares:
% \begin{ClozeArea}[<columns>]{<rows>} ... \end{ClozeArea}
\NewDocumentEnvironment{ClozeArea}{O{34} m +b}{%
  \par\medskip\noindent
  \pgfmathsetlengthmacro{\ClozeAreaWidth}{#1*5mm}%
  \pgfmathsetlengthmacro{\ClozeAreaHeight}{#2*5mm}%
  \hfill
  \IfClozeTF{%
    \begin{tikzpicture}
      \path[use as bounding box]
        (0,0) rectangle (\ClozeAreaWidth,\ClozeAreaHeight);
      \draw[step=5mm,black!14,line width=0.2pt]
        (0,0) grid (\ClozeAreaWidth,\ClozeAreaHeight);
      \draw[black!25,line width=0.3pt]
        (0,0) rectangle (\ClozeAreaWidth,\ClozeAreaHeight);
    \end{tikzpicture}%
  }{%
    \begin{minipage}{\ClozeAreaWidth}#3\end{minipage}%
  }%
  \hfill\mbox{}%
  \par\medskip
}{}

% Shared geometry for branching and merging equations.
\newlength{\EquationTemplateUpperWidth}
\newlength{\EquationTemplateLowerWidth}
\newlength{\EquationTemplateMaxWidth}
\NewDocumentCommand{\EquationTemplateMeasure}{m m}{%
  \settowidth{\EquationTemplateUpperWidth}{\ensuremath{\displaystyle #1}}%
  \settowidth{\EquationTemplateLowerWidth}{\ensuremath{\displaystyle #2}}%
  \ifdim\EquationTemplateUpperWidth>\EquationTemplateLowerWidth
    \setlength{\EquationTemplateMaxWidth}{\EquationTemplateUpperWidth}%
  \else
    \setlength{\EquationTemplateMaxWidth}{\EquationTemplateLowerWidth}%
  \fi
}
\NewDocumentCommand{\EquationTemplateOverlay}{m m m}{%
  \makebox[0pt][l]{%
    \hspace{#2}\raisebox{#1}[0pt][0pt]{#3}}%
}
\NewDocumentCommand{\EquationTemplateEquals}{m}{%
  \rotatebox{#1}{\ensuremath{\displaystyle =}}%
}
\NewDocumentCommand{\EquationTemplateUse}{m}{%
  \ifmmode
    \vcenter{\hbox{#1}}%
  \else
    #1%
  \fi
}
\NewDocumentCommand{\MathHighlight}{m}{%
  \begingroup
  \setlength{\fboxsep}{1.5pt}%
  \colorbox{yellow!65}{\ensuremath{\displaystyle #1}}%
  \endgroup
}

% One expression continued in two alternative directions:
% \EquationBranches{<start>}{<upper continuation>}{<upper note>}
%                  {<lower continuation>}{<lower note>}
\NewDocumentCommand{\EquationBranchesContent}{m m m m m}{%
  \begingroup
  \EquationTemplateMeasure{#2}{#4}%
  \ensuremath{\displaystyle #1}%
  \hspace{0.15cm}%
  \EquationTemplateOverlay{0.7em}{0pt}{\EquationTemplateEquals{35}}%
  \EquationTemplateOverlay{1.4em}{0.55cm}{%
    \makebox[\EquationTemplateMaxWidth][l]{\ensuremath{\displaystyle #2}}%
    \hspace{0.55cm}#3}%
  \EquationTemplateOverlay{-0.7em}{0pt}{\EquationTemplateEquals{-35}}%
  \EquationTemplateOverlay{-1.4em}{0.55cm}{%
    \makebox[\EquationTemplateMaxWidth][l]{\ensuremath{\displaystyle #4}}%
    \hspace{0.55cm}#5}%
  \rule[-1.8em]{0pt}{4.3em}%
  \endgroup
}

\NewDocumentCommand{\EquationBranches}{m m m m m}{%
  \EquationTemplateUse{%
    \EquationBranchesContent{#1}{#2}{#3}{#4}{#5}}%
}

% Two alternative expressions leading to one common result:
% \EquationMerge{<upper note>}{<upper expression>}{<lower note>}
%                {<lower expression>}{<result>}
\NewDocumentCommand{\EquationMergeContent}{m m m m m}{%
  \begingroup
  \EquationTemplateMeasure{#2}{#4}%
  \makebox[\EquationTemplateMaxWidth][l]{%
    \EquationTemplateOverlay{1.4em}{0pt}{%
      \makebox[0pt][r]{#1\hspace{0.55cm}}%
      \makebox[\EquationTemplateMaxWidth][r]{\ensuremath{\displaystyle #2}}}%
    \EquationTemplateOverlay{-1.4em}{0pt}{%
      \makebox[0pt][r]{#3\hspace{0.55cm}}%
      \makebox[\EquationTemplateMaxWidth][r]{\ensuremath{\displaystyle #4}}}%
    \rule[-1.8em]{0pt}{4.3em}}%
  \hspace{0.15cm}%
  \EquationTemplateOverlay{0.7em}{0pt}{\EquationTemplateEquals{-35}}%
  \EquationTemplateOverlay{-0.7em}{0pt}{\EquationTemplateEquals{35}}%
  \hspace{0.55cm}%
  \ensuremath{\displaystyle #5}%
  \endgroup
}

\NewDocumentCommand{\EquationMerge}{m m m m m}{%
  \EquationTemplateUse{%
    \EquationMergeContent{#1}{#2}{#3}{#4}{#5}}%
}

\newcommand{\R}{\mathbb{R}}
\newcommand{\N}{\mathbb{N}}
\newcommand{\Q}{\mathbb{Q}}
\newcommand{\Z}{\mathbb{Z}}
\newcommand{\C}{\mathbb{C}}
\DeclareMathOperator{\grad}{grad}
\DeclareMathOperator{\arsinh}{arsinh}
\DeclareMathOperator{\arcosh}{arcosh}
\DeclareMathOperator{\artanh}{artanh}
\newcommand{\corresponds}{\mathrel{\widehat{=}}}
